% Entwurf 4.1 (2026-10-17), Claude entwirft, der Nutzer redigiert.
% Schlüssel: Stommel, Stojek, Korte 2018 steht in der BA-Bib unter Sto18 (Sto18 ist dort
% Stoffregen, Motorradtechnik); hier wird Sto18 zitiert, damit der Eintrag auflöst.
\section{Netzkonvergenz und Elementtyp}
\label{sec:netzkonvergenz}

Vor dem Einfrieren des Radnetzes wurden Netzdichte und Elementformulierung geprüft. Einer unzureichenden Diskretisierung begegnet nach Stommel, Stojek und Korte die h-Methode, das Netz wird verfeinert, bis sich die Ergebnisse benachbarter Stufen nur noch wenig ändern \cite[S.~424]{Sto18}. Die Konvergenzstudie läuft vor der eigentlichen Auswertung an einem kleinen Modell unter vergleichbarer Last mit mehreren Elementtypen gegen eine Referenz \cite[S.~433]{Sto18}. Die Elementgröße wird dabei so groß gewählt, wie die Beanspruchungsgradienten es zulassen \cite[S.~465]{Sto18}. Als kleines Modell dient der Stirnschnitt des Radsektors im ebenen Dehnungszustand, weil das Volumennetz ein Auszug dieses Schnitts über die Zahnbreite ist. Ein eigener ebener Löser rechnet das Vierknotenelement selektiv reduziert integriert wie CPE4 und C3D8 sowie mit inkompatiblen Moden wie CPE4I und C3D8I. \Note{PH: Quelle Abaqus-Dokumentation} Am Einzahnsektor mit 859 Knoten stimmt er in der Abaqus Learning Edition mit CPE4 auf $7 \cdot 10^{-5}$ überein, die Vollversion bestätigte jede Stufe der Leiter auf $0{,}01$\,MPa.

\figref{fig:lastfall_eben} zeigt den Lastfall. Die Last von $16$\,Nm am Rad wirkt am äußeren Einzeleingriffspunkt B mit dem längsten Hebelarm auf einen Zahn, als Linienlast von $43{,}66$\,N/mm. Sie ist als Hertzsche Pressungsverteilung nach \cite{Nie19} mit der halben Druckbreite $0{,}322$\,mm und der größten Pressung $86{,}4$\,MPa normal zur Flanke aufgebracht. Der Werkstoff ist isotrop mit dem Elastizitätsmodul $E = 2282$\,MPa bei $80$\,°C und der Querkontraktionszahl $0{,}30$, Bohrung und Schnittebenen sind fest. Bewertet werden die größte Hauptspannung $\sigma_{\mathrm{1}}$ an der Fußoberfläche der belasteten Flanke als Knotenmittel, die Verschiebung der Kopfmitte und die Formänderungsenergie des Sektors. Das Netz der Vorlage hat 20 Elemente über der Zahnhöhe, 80 je Zahnlücke am Fuß und 10 über der Zahndicke. Verfeinert wurde gleichmäßig bis auf ein Viertel der Elementgröße und jede der drei Richtungen für sich. Der Grenzwert für die Elementgröße null folgt nach Richardson aus den letzten drei Stufen mit dem Ansatz
\begin{equation}
f(h) = f_{\mathrm{0}} + C\,h^{p},
\label{eq:richardson}
\end{equation}
wobei der Exponent $p$ aus den Stufen mitbestimmt wird. \Note{PH: Formelzeichen f, f0, C, h und p in die Nomenklatur aufnehmen; die Quellen Bathe 1996, Gl. (4.101), und Dahmen und Reusken 2022, Gl. (10.35), fehlen im Literaturverzeichnis}

\begin{figure}[htbp]
\centering
\includegraphics[width=0.8\textwidth]{kap4/lastfall_eben}
\caption{Lastfall der Konvergenzstudie am ebenen Radsektor: Pressungsverteilung am Punkt B bei $16$\,Nm am Rad und größte Hauptspannung je Element im Netz der Vorlage.}
\label{fig:lastfall_eben}
\end{figure}

\figref{fig:konvergenz_fussspannung} und \tabref{tab:netzleiter} zeigen die Fußspannung. Im Netz der Vorlage beträgt sie $125{,}4$\,MPa. Eine Verfeinerung über der Zahnhöhe oder entlang der Fußrundung ändert sie um weniger als $0{,}1$\,\%. Allein die Richtung normal zur Fußoberfläche, die Sehnen der Zahndicke, hebt sie auf $135{,}1$, $138{,}9$ und $140{,}9$\,MPa. Diese Reihe konvergiert mit der Ordnung $0{,}85$ gegen $148{,}1$\,MPa, die Vorlage liegt also $15{,}4$\,\% darunter, die doppelte Zahndicke noch $8{,}8$\,\% und die vierfache $4{,}9$\,\%. Die gleichmäßige Reihe gibt über drei Stufen $150{,}0$\,MPa, mit einer vierten Stufe von 55\,193 Knoten $148{,}2$\,MPa. Der langsame Verlauf liegt an der Formulierung, das Element hält den volumetrischen Dehnungsanteil je Element konstant, während der hydrostatische Anteil an der Kerboberfläche von null aus steil nach innen ansteigt. Die Breite der Pressungsverteilung verschiebt die Absolutwerte um etwa $5$\,\%, das Konvergenzverhalten der Reihen nicht.

\begin{figure}[htbp]
\centering
\includegraphics[width=0.8\textwidth]{kap4/konvergenz_fussspannung}
\caption{Größte Hauptspannung an der Fußoberfläche über der relativen Elementgröße für die gleichmäßige Verfeinerung und die drei Richtungen des Netzes, Formulierung von CPE4, mit den Grenzwerten nach \eqref{eq:richardson}.}
\label{fig:konvergenz_fussspannung}
\end{figure}
% Bildwunsch: Achsenbeschriftung "1 = Netz der Vorlage" statt "gestriges Netz"; die gleichmäßige
% Reihe mit der vierten Stufe (55 193 Knoten), damit Bild und Text denselben Grenzwert nennen.
% Bildwunsch: eine zweite Fassung der Leiter Zahndicke mit CPE4R, CPE4 und CPE4I aus den
% Abaqus-Decks (Tabelle unten), die die Unterschrift "drei Elementformulierungen" aus
% bilder.yaml einlöst; das heutige Bild zeigt nur die Formulierung von CPE4.

Für die Verformung am Zahnkopf genügt die Dichte der Vorlage. Die Verschiebung der Kopfmitte beträgt dort $203{,}8$\,\textmu m und ändert sich in keiner Reihe um mehr als $0{,}2$\,\%, die Reihe Zahndicke läuft mit der Ordnung $4{,}25$ gegen $203{,}47$\,\textmu m. Die Formänderungsenergie des Sektors ändert sich über alle Stufen um höchstens $0{,}14$\,\%, und alle drei Elementformulierungen liegen in der Kopfverschiebung innerhalb von $1{,}2$\,\% beieinander.

% Bild verfügbar, wegen des Seitenbudgets nicht eingebunden (Vorschlag: Anhang E):
% \begin{figure}[htbp]
% \centering
% \includegraphics[width=0.8\textwidth]{kap4/konvergenz_kopfverschiebung}
% \caption{Verschiebung der Kopfmitte über der relativen Elementgröße für die gleichmäßige Verfeinerung und die drei Richtungen des Netzes, Formulierung von CPE4.}
% \label{fig:konvergenz_kopfverschiebung}
% \end{figure}

\begin{table}[htbp]
\centering
\caption{Netzleiter über der Zahndicke am ebenen Radsektor, gerechnet auf der Abaqus-Vollversion: größte Hauptspannung an der Fußoberfläche und Verschiebung der Kopfmitte je Elementformulierung, Grenzwert und Ordnung nach \eqref{eq:richardson}.}
\label{tab:netzleiter}
\begin{tabular}{lllllll}
\hline
Elemente über der & Knoten je & \multicolumn{3}{l}{$\sigma_{\mathrm{1}}$ am Fuß in MPa} & \multicolumn{2}{l}{Kopfmitte in \textmu m} \\
Zahndicke & Schicht & CPE4R & CPE4 & CPE4I & CPE4 & CPE4I \\
\hline
10 & 3\,719 & $114{,}85$ & $125{,}42$ & $141{,}60$ & $203{,}8$ & $202{,}0$ \\
20 & 7\,021 & $127{,}94$ & $135{,}13$ & $147{,}61$ & $203{,}5$ & $203{,}1$ \\
30 & 10\,323 & $133{,}57$ & $138{,}91$ & $148{,}50$ & $203{,}5$ & $203{,}3$ \\
40 & 13\,625 & $136{,}69$ & $140{,}90$ & $148{,}62$ & $203{,}5$ & $203{,}4$ \\
\hline
Grenzwert & & $150{,}5$ & $148{,}1$ & $148{,}7$ & $203{,}5$ & $203{,}6$ \\
Ordnung & & $0{,}71$ & $0{,}85$ & $4{,}63$ & $4{,}25$ & $1{,}8$ \\
\hline
\end{tabular}
\end{table}

Die Elementformulierung wiegt so schwer wie die Dichte (vgl. \tabref{tab:netzleiter}). Mit reduzierter Integration, dem ebenen Gegenstück des C3D8R im Referenznetz, liegt der Knotenwert im Netz der Vorlage bei $114{,}9$\,MPa und damit rund $23$\,\% unter dem Grenzwert. Die Hourglass-Energie bleibt dabei unter $0{,}04$\,\% der Formänderungsenergie und ist nach der Energiebilanz unkritisch \cite[S.~428]{Sto18}. Mit inkompatiblen Moden liefert dasselbe Netz $141{,}6$\,MPa und damit $4{,}7$\,\% unter dem Grenzwert, die doppelte Zahndicke $147{,}6$\,MPa mit $0{,}7$\,\%, und die Reihe konvergiert mit der Ordnung $4{,}6$. Alle drei Formulierungen weisen auf denselben Grenzwert zwischen 148 und $149$\,MPa. Für die Wahl der Dichte wurde eine Leiter aus Zahndicke und Fußschichten mit diesem Element gerechnet und nach \eqref{eq:richardson} über alle Stufen ausgeglichen. Die Toleranz betrug $2$\,\% für die Fußspannung und $1$\,\% für die Kopfverschiebung. Die Stufe mit 14 Elementen über der Zahndicke und 4 Fußschichten liegt $1{,}8$\,\% unter dem Grenzwert, die Stufe mit 20 und 6 noch $0{,}7$\,\%. Die übrigen Anzahlen wurden einzeln vergröbert, solange beide Größen in der Toleranz blieben, nur die Zahnhöhe behielt ihre 18 Elemente, weil der Kontakt drei Elemente je Hertzscher Halbbreite verlangt. Das empfohlene Netz hat 18 Elemente über der Zahnhöhe und eines am Kopfkantenbruch, 60 am Zahnfuß, 14 über der Zahndicke mit 4 Fußschichten und 4 Kranzringe. Mit 3\,890 Knoten je Schicht liegt seine Fußspannung $1{,}9$\,\% unter dem Grenzwert.

Die ersten Stellungsdateien mit C3D8I und geometrisch nichtlinearer Rechnung brachen im Schritt mit $12$\,Nm ab. Warnungen negativer Eigenwerte der Systemmatrix erschienen ab dem ersten Lastinkrement und bei allen drei Kontaktformulierungen gleich. Ein Elementtest am Grobmodell der Learning Edition mit einem Zahn, zwei Schichten und 414 Knoten trennte die Ursache. Mit C3D8I und dem Schalter NLGEOM=YES erscheinen 29 Warnungen ab dem Schritt mit $8$\,Nm, Inkrement 1, während dasselbe Modell mit C3D8 oder C3D8R geometrisch nichtlinear ohne eine einzige Warnung rechnet. Die Elementdokumentation von Abaqus beschränkt Elemente mit inkompatiblen Moden auf kleine Druckdehnungen. \Note{PH: Quelle Abaqus-Dokumentation} Am Pilotnetz mit 20 Schichten und 81\,690 Knoten laufen beide C3D8-Dateien alle vier Schritte ohne negativen Eigenwert durch, in 140 und $161$\,s gegen $173$\,s für C3D8I. Die geometrische Nichtlinearität ändert bis $16$\,Nm das Stützmoment um $0{,}1$\,\%, die Ritzeldrehung um $0{,}5$\,\% und die größte Pressung um $0{,}7$\,\%. Der Wechsel von C3D8I auf C3D8 ändert diese Größen um weniger als $1$\,\%.

Daraus folgt die Regel je Werkstoffstufe. Die elastischen Stufen W1 und W3 werden geometrisch linear mit C3D8I gerechnet, weil ein lineares Problem linear gerechnet wird \cite[S.~474]{Sto18} und die Fußspannung dort innerhalb von $2$\,\% am Grenzwert liegt. Die elastisch-plastischen Stufen W2 und W4 brauchen die geometrisch nichtlineare Rechnung, da sich Kunststoffe schon bei geringen Lasten plastisch verformen und die geometrische Nichtlinearität mit der Dehnung wächst \cite[S.~2, 52, 474]{Sto18}. Sie werden mit C3D8 auf demselben Netz gerechnet, Elementnummern und Faserorientierung bleiben erhalten. Die Leiter mit der Formulierung von C3D8 konvergiert an der Fußoberfläche nur mit der Ordnung $0{,}81$ und erreicht auf keiner Stufe $2$\,\%. Bei 14 Elementen und 4 Fußschichten liegt sie $12{,}0$\,\% unter dem Grenzwert, bei 40 und 12 mit 286\,000 Knoten im Volumen noch $4{,}9$\,\%. Fußspannungen einer C3D8-Stufe sind deshalb mit denen einer C3D8I-Stufe nur auf gleichem Elementtyp oder gegen den extrapolierten Grenzwert vergleichbar.

Die Zahl der Schichten über der Zahnbreite prüfte das ebene Modell nicht. Volumendecks der Vorlagendichte mit 20, 40 und 80 Schichten, gerechnet auf der Vollversion mit C3D8I, geben in der Radmitte $204{,}87$, $204{,}93$ und $204{,}95$\,\textmu m sowie $143{,}49$, $143{,}53$ und $143{,}55$\,MPa. Nur der Wert unmittelbar an der Stirnfläche fällt mit feinerer Teilung weiter von $108{,}3$ auf $96{,}8$\,MPa, er gehört zum freien Rand und ist keine Zielgröße. Gegenüber dem Volumenmodell nimmt der ebene Dehnungszustand das Rad in der Radmitte um $1{,}4$\,\% zu steif. Die Pilotläufe der Stellungsdatei mit 40 und 80 Schichten ließen Momente und Drehung auf vier Stellen unverändert und hoben die größte Pressung um $1$\,\%. Eingefroren wurde deshalb das Netz mit 20 Schichten in der empfohlenen Dichte, mit fünf Kranzringen bis zur Bohrung von $33$\,mm, dem Innendurchmesser des Kunststoffbereichs am realen Rad. Der Stahleinleger darin ist als feste Fesselung abgebildet. Das eingefrorene Netz mit Ringkörper hat 82\,425 Knoten, der Taschenkörper folgt in \secref{sec:stufen_raten_koerper}, die Tabellen aller Reihen stehen in \secref{sec:anhang_e}.

% Bild verfügbar, wegen des Seitenbudgets nicht eingebunden (Vorschlag: Anhang E):
% \begin{figure}[htbp]
% \centering
% \includegraphics[width=0.8\textwidth]{kap4/fussspannung_verlauf_eben}
% \caption{Verlauf der größten Hauptspannung entlang der Fußrundung der belasteten Flanke für die gleichmäßigen Stufen 1 bis 3 im ebenen Modell.}
% \label{fig:fussspannung_verlauf_eben}
% \end{figure}

% Herkunft:
% Absatz 1 (Vorgehen, Löser, Prüfung): 20_code/00_development_documentation/fe_mesh_convergence.md
%   Abschnitt 2 (Stommel 7.7.1 S. 424, 7.7.3 S. 433, 8.2.2 S. 465; ebener Sektor als Auszugsnetz;
%   Löser; Prüfung Learning Edition 7e-5) und Abschnitt 5a (Vollversion: CPE4 alle gedruckten
%   Stellen, CPE4I 0,01 MPa und 0,01 µm); 80_output/fe/kst_e/convergence/verification.md
%   (859 Knoten, Tabelle); architecture_decisions.md ADR-116 Punkt 12.
% Absatz 2 (Lastfall, Größen, Reihen, Formel): fe_mesh_convergence.md Abschnitt 2 (Last, Hertz nach
%   Nie19 Tab. 13.1, E = 2282 MPa, ν = 0,30, Fesselung, bewertete Größen, Reihen, Richardson nach
%   Bathe und Dahmen/Reusken) und Abschnitt 1 Tabelle (Vorlage 20 / 80 / 10);
%   Bild load_case.png (F' = 43,66 N/mm, b_H = 0,322 mm, p_0 = 86,4 MPa).
% Absatz 3 (Fußspannung): fe_mesh_convergence.md Abschnitt 3 Tabelle und Extrapolation (Zahndicke
%   p = 0,85 gegen 148,1; Fehler −15,4 / −8,8 / −6,2 / −4,9 %; gleichmäßig vier Stufen p = 0,84
%   gegen 148,2 mit 55 193 Knoten; Lastbreite etwa 5 %), Abschnitt 4 Punkt 2 (Ursache, erste
%   Schicht 0,02 mm); 80_output/fe/kst_e/bilder_V2/convergence/convergence.md (gleichmäßig drei
%   Stufen p = 0,72 gegen 150,0; Zahnhöhe und Zahnfuß ≤ 0,03 %); Bild convergence_root_stress.png.
% Absatz 4 (Kopfverschiebung, Energie): bilder_V2/convergence/convergence.md Extrapolation Reihe
%   Zahndicke (u Kopfmitte p = 4,25 gegen 203,465 µm); fe_mesh_convergence.md Abschnitt 3
%   (≤ 0,2 %, U ≤ 0,14 %) und Abschnitt 4 Punkt 1 (drei Formulierungen innerhalb 1,2 %).
% Tabelle: 80_output/fe/kst_e/full_licence_result/comparison.md (ebene Decks base/t2/t3/t4 je
%   CPE4R/CPE4/CPE4I; Grenzwerte 150,47 (0,71), 148,13 (0,85), 148,66 (4,63)); Kopfmitte CPE4
%   Grenzwert und Ordnung aus bilder_V2/convergence/convergence.md (203,465, p = 4,25), CPE4I
%   Grenzwert 203,6 µm (p = 1,8) aus fe_mesh_convergence.md Abschnitt 3 (eigener Löser).
% Absatz 5 (Elementtyp, Empfehlung): fe_mesh_convergence.md Abschnitt 4 Punkte 3 und 4, Abschnitt
%   5a (Hourglass < 0,04 %), Abschnitt 5b (Leiter 14/4 −1,82 %, 20/6 −0,72 %; Reduktion; Zahnhöhe
%   Untergrenze 19; Empfehlung 18 + 1 / 60 / 14 / 4 / 4 / 2, 3 890 Knoten, −1,9 %);
%   80_output/fe/kst_e/recommendation/recommendation.md; ADR-116 Punkte 13 bis 15.
% Absatz 6 (Nichtlinearität, Elementtest, Bestätigung): fe_mesh_convergence.md Abschnitt 5c;
%   80_output/fe/kst_e/le_tests/element_nlgeom/ELEMENT_TEST.md (414 Knoten, 29 Warnungen ab 8 Nm
%   Inkrement 1; das Tagebuch fe_simulation_log.md zählt für denselben Lauf 30 negative
%   Eigenwerte); ADR-116 Punkte 18 und 20 (140 s / 161 s / 173 s; das Tagebuch nennt für
%   pilot_20layers_c3d8_variants_gerechnet 157 s und 174 s Wandzeit); known_limits.md FE-14.
% Absatz 7 (Regel je Stufe): ADR-116 Punkte 19 (b), 19 (c) und 21 (Stommel S. 2, 52, 474);
%   fe_mesh_convergence.md Abschnitt 5c Tabelle (C3D8-Leiter Ordnung 0,81, −12,0 / −8,8 / −4,9 %,
%   286 125 Knoten bei 20 Schichten); known_limits.md FE-10.
% Absatz 8 (Schichten, Bohrung, Einfrieren): full_licence_result/comparison.md Volumendecks
%   (C3D8I 20/40/80 Schichten); fe_mesh_convergence.md Abschnitt 5a (Stirnfläche 108 → 97 MPa,
%   ebener Zustand 1,4 % zu steif) und Abschnitt 6; ADR-116 Punkt 22 (Piloten 40/80 Schichten,
%   Bohrung 16,5 mm Radius = 33 mm, Stahleinleger als Fesselung, 82 425 Knoten);
%   known_limits.md FE-08 und FE-09.
